\section{Zero-locus numerics as verification}\label{sec:verification}

The preceding sections establish the complete spectral geometry of the quartic transfer problem by exact algebraic and analytic methods.
Numerical computation of the zeros of \(Z_m(y)\) and of the dominant-modulus competition serves here as an independent consistency check, not as primary evidence.

\subsection{Validation targets from the exact theory}

Set
\[
\Omega:=\CC\setminus\bigl(\Branch\cup\{-1,0,1\}\bigr).
\]
For each fixed $y \in \Omega$, the spectral polynomial $\Pi(\cdot, y)$ has four simple roots $\lambda_1(y), \dots, \lambda_4(y)$, and Appendix~\ref{app:vacuum-edge} gives the spectral decomposition
\[
Z_m(y)=\sum_{j=1}^4 C_j(y)\lambda_j(y)^m,
\qquad
C_j(y)=\frac{\lambda_j(y)^3\,N(1/\lambda_j(y),y)}{\partial_\lambda\Pi(\lambda_j(y),y)}.
\]
The Beraha--Kahane--Weiss theorem, in the form used here \cite{BerahaKahaneWeiss1975,Sokal2004}, gives two alternatives for a limit point of zeros of \(\{Z_m(y)\}_{m\ge0}\): either two or more eigenvalues of maximal modulus tie, or a unique dominant eigenvalue exists but its coefficient vanishes.
On \(\Omega\), the second alternative does not occur.
Indeed, for a simple root \(\lambda_j(y)\) one has \(\partial_\lambda\Pi(\lambda_j(y),y)\neq 0\), and Theorem~\ref{thm:minimal-order} shows that \(N(t,y)\) and \(D(t,y)\) are coprime for \(y\notin\{-1,0,1\}\).
Since \(D(1/\lambda_j(y),y)=0\), it follows that \(N(1/\lambda_j(y),y)\neq 0\), hence \(C_j(y)\neq 0\) on \(\Omega\).
Therefore, on the generic set \(\Omega\), the limiting accumulation set is governed solely by dominant equimodularity:
\[
\bigl\{y \in \Omega : \exists\, i \neq j \text{ with }
\abs{\lambda_i(y)} = \abs{\lambda_j(y)} = r(y)\bigr\}.
\]
Branch points in $\Branch$ arise as natural endpoints or singularities of such equimodular curves.

\subsection{A numerical consistency figure}

Figure~\ref{fig:validation} summarizes a representative numerical check generated directly from the order-$4$ recurrence \eqref{eq:recurrence} and the quartic spectral polynomial \eqref{eq:spectral-poly}.
No fitting parameters are introduced: all markers correspond to the exact branch values from Theorem~\ref{thm:elliptic-reduction} and Theorem~\ref{thm:ep-classification}.

\begin{figure}[t]
\centering
\IfFileExists{figures/zero_locus_validation.pdf}{%
\includegraphics[width=\textwidth]{figures/zero_locus_validation.pdf}%
}{%
\fbox{\parbox{\dimexpr\textwidth-2\fboxsep-2\fboxrule}{%
\centering\vspace{2cm}%
[Validation figure not included in this build.]%
\vspace{2cm}}}%
}
\caption{
\textbf{Left:} ordered root moduli of $\Pi(\lambda,y)=0$ along the real axis.
At $y=0$, the top two moduli meet, confirming that the dominant sheet over $y>0$ terminates there.
At $y=1$, only the lower pair collides while the dominant root remains isolated.
At $y=\ySub$, the colliding double root lies strictly below the dominant modulus, as predicted by Theorem~\ref{thm:ep-classification}.
\textbf{Center:} zeros of $Z_{80}(y)$ and $Z_{120}(y)$ in the complex $y$-plane.
The finite-volume zero clouds organize around the same branch-value window and terminate at $y=0$ rather than at $y=\ySub$.
The algebraically forced multiplicity-two contribution at $y=-1$ is pinned for every $m\ge 2$, while no comparable zero accumulation appears near $y=1$.
\textbf{Right:} scaled negative zeros $-m^2y_{m,k}$ for $k=1,\dots,4$.
Filled markers are exact zeros of $Z_m$ for $m=80,120,200$; open markers are the two-term asymptotic prediction from Theorem~\ref{thm:vacuum-zero-quantization}; the black line is the leading odd-integer law $\pi^2(2k-1)^2/2$.
The data verify both the $m^{-2}$ edge scale and the odd-integer quantization inherited from the square-root exceptional point.
}
\label{fig:validation}
\end{figure}

The left panel is the most direct numerical check of the branch classification.
Sorting the four spectral roots by modulus for real $y$ shows that the dominant-modulus collision occurs only at $y = 0$.
The collision at $y = 1$ is visibly subdominant, and the collision at $y = \ySub$ occurs below the dominant complex-conjugate pair.
This is the numerical counterpart of the exact inequalities proved in Section~\ref{sec:ep-classification}.

The center panel gives the finite-volume zero-cloud counterpart of this classification.
For $m=80$ and $m=120$, the zeros accumulate in the branch-value window attached to the real-axis competition picture, end at the square-root endpoint $y=0$, and retain the guaranteed order-at-least-two boundary zero at $y=-1$.
The absence of comparable zero concentration near $y = 1$ is consistent with the fact that $y=1$ is not a branch point of the dominant sheet over $y>0$.

The right panel checks the local dominant-EP2 endpoint law specialized in Theorem~\ref{thm:vacuum-zero-quantization}.
For the first several negative zeros, the scaled values $-m^2y_{m,k}$ line up with the odd-integer sequence $\pi^2(2k-1)^2/2$ and are already well captured by the explicit $m^{-3}$ correction in Theorem~\ref{thm:vacuum-zero-quantization}.
This is the finite-volume numerical counterpart of the cosine normal form \eqref{eq:vacuum-scaling} and of the counting law from Corollary~\ref{cor:edge-zero-counting}.

\medskip
\noindent\textbf{Reproducibility.}
The numerical conventions used for Figure~\ref{fig:validation} are recorded explicitly in Appendix~\ref{app:pseudocode}.
Concretely, the left panel samples \(1500\) equally spaced real points on \([-1.6,1.2]\) and orders the four roots of \(\Pi(\lambda,y)\) by decreasing modulus.
The center panel plots all roots of \(Z_{80}\) and \(Z_{120}\) computed via \texttt{numpy.roots}, with the guaranteed multiplicity-two contribution at \(y=-1\) inserted explicitly after factoring out \((y+1)^2\).
The right panel uses \(m=80,120,200\) and the first four negative real zeros nearest \(0\), selected by the criterion \(\abs{\Im y}<10^{-8}\) together with \(-1<\Re y<0\).
All computations use IEEE double precision, and the real branch value \(y_{\mathrm{sub}}\) is extracted from \(256y^3+411y^2+165y+32\) by the filter \(\abs{\Im y}<10^{-10}\).
